\chapter{Glossaries}

OmniLaTeX uses the \texttt{glossaries-extra} package to provide glossary,
acronym, symbol, and index functionality.  All glossary infrastructure is
loaded by \texttt{omnilatex-glossary.sty}; you do \emph{not} need to load
\texttt{glossaries} or \texttt{glossaries-extra} yourself.

\section{glossaries-extra Setup}

The package is initialised with the \texttt{nomain} option so that
\texttt{omnilatex-glossary.sty} retains full control over which glossary
types are created.  The following glossary types are registered:

\begin{itemize}
    \item \texttt{acronyms} -- abbreviations and acronyms
    \item \texttt{symbols} -- mathematical and physical symbols
    \item \texttt{index} -- back-of-book index entries
    \item \texttt{numbers} -- physical constants and named numbers
\end{itemize}

Each type receives an appropriate sorting and labelling strategy via
\texttt{glossaries-extra}'s \texttt{sort} and \texttt{label} keys.

\section{Custom Glossary Fields}

OmniLaTeX extends the default entry keys with several custom fields that
are recognised by the symbol and constant glossary styles.

\begin{table}[htbp]
    \centering
    \caption{Custom glossary fields added by OmniLaTeX}
    \label{tab:glossary-fields}
    \begin{tabular}{@{} l l p{6cm} @{}}
        \toprule
        \textbf{Field} & \textbf{Type} & \textbf{Purpose} \\
        \midrule
        unit                  & text & SI unit associated with a symbol \\
        specific              & text & Specific context or domain qualifier \\
        firstname             & text & Person's first name (for named entries) \\
        value                 & text & Numerical value of a constant \\
        international-symbol  & text & International symbol variant \\
        alternative-symbol    & text & Alternative or legacy symbol \\
        \bottomrule
    \end{tabular}
\end{table}

These fields are accessed inside glossary styles with
\texttt{\textbackslash glsuseri}, \texttt{\textbackslash glsuserii}, etc.

\section{Shortcut Commands}

OmniLaTeX provides shorthand macros that create and reference glossary
entries in a single call.  Each command targets a specific glossary type.

\begin{table}[htbp]
    \centering
    \caption{Glossary shortcut commands}
    \label{tab:glossary-shortcuts}
    \begin{tabular}{@{} l l l @{}}
        \toprule
        \textbf{Command} & \textbf{Glossary Type} & \textbf{Example} \\
        \midrule
        \verb|\sym{Key}|   & symbol       & \verb|\sym{rho}|     \\
        \verb|\sub{Key}|   & subscript    & \verb|\sub{max}|     \\
        \verb|\name{Key}|  & name         & \verb|\name{Einstein}| \\
        \verb|\cons{Key}|  & constant     & \verb|\cons{g}|      \\
        \verb|\idx{Key}|   & index        & \verb|\idx{entropy}| \\
        \verb|\abb{Key}|   & abbreviation & \verb|\abb{CPU}|     \\
        \verb|\num{Key}|   & number       & \verb|\num{pi}|      \\
        \bottomrule
    \end{tabular}
\end{table}

The first invocation of each command creates the entry (with sensible
defaults); subsequent calls expand to the formatted representation.

\section{Symbol Glossary}

The \texttt{symbols} glossary uses the \texttt{symbunitlong} style, which
displays each symbol together with its unit in a long-table layout.
Entries that carry a \texttt{specific} field are annotated with a
footnote-style marker so the reader can distinguish homographic symbols
(e.g.\ $c$ for speed of light vs.\ $c$ for specific heat capacity).

\section{Constants Glossary}

The \texttt{numbers} glossary employs the \texttt{numberlong} style,
which prints each constant's name, symbol, and numerical value side by
side.  The \texttt{value} field is rendered in a fixed-width column so
that mantissas align vertically.

\section{Printing}

Glossaries are printed with \texttt{\textbackslash printunsrtglossary}.
Two calls are typically placed in the back matter:

\begin{verbatim}
\printunsrtglossary[type=symbols, style=symbunitlong,
                    title={Nomenclature}]
\printunsrtglossary[type=acronyms,
                    style=\OmniLaTeXAcronymStyle]
\end{verbatim}

\texttt{\textbackslash OmniLaTeXAcronymStyle} is a custom style defined
in \texttt{omnilatex-glossary.sty} that formats acronyms with their long
form on first use and the short form thereafter.

\section{Index}

The back-of-book index is generated with

\begin{verbatim}
\printunsrtindex[style=bookindex]
\end{verbatim}

Index entries are created via \verb|\idx{key}|, which delegates to
\texttt{\textbackslash gls} with the \texttt{index} glossary type.
